\clearpage
\begin{explanationbox}
\paragraph{一、先看具体曲线：贴合发生在零点附近}
\begin{center}
\begin{tikzpicture}[font=\small]
\draw[black!30,rounded corners=2pt] (0,0) rectangle (14.4,8.2);
\node[anchor=west,font=\small\bfseries] at (0.5,7.72) {指数：先在零点附近比较};
\begin{axis}[width=11.2cm,height=6.1cm,at={(2.05cm,1.25cm)},anchor=south west,
axis lines=left,xlabel={$x$},tick label style={font=\small},label style={font=\small},
legend style={font=\small,draw=none,fill=white,at={(0.03,0.97)},anchor=north west},
legend cell align=left,unbounded coords=jump,xmin=-0.8,xmax=0.8,ymin=0,ymax=2.6,ylabel={$y$}]
\addplot[blue,very thick] table[x=x,y=exp] {assets/tikz/F003_samples.dat};
\addlegendentry{$e^x$}
\addplot[orange!85!black,thick,dashed] table[x=x,y=re] {assets/tikz/F003_samples.dat};
\addlegendentry{$R_E(x)$}

\end{axis}
\end{tikzpicture}
\end{center}

蓝线是 $e^x$，橙虚线是 $R_E$。两线近乎重合，只能说明在当前尺度下差距很小，不能说明相等。
\begin{center}
\begin{tikzpicture}[font=\small]
\draw[black!30,rounded corners=2pt] (0,0) rectangle (14.4,8.2);
\node[anchor=west,font=\small\bfseries] at (0.5,7.72) {指数：把两条曲线的微小差距放大来看};
\begin{axis}[width=11.2cm,height=6.1cm,at={(2.05cm,1.25cm)},anchor=south west,
axis lines=left,xlabel={$x$},tick label style={font=\small},label style={font=\small},
legend style={font=\small,draw=none,fill=white,at={(0.03,0.97)},anchor=north west},
legend cell align=left,unbounded coords=jump,ymode=log,xmin=0.1,xmax=1,ymin=1e-9,ymax=0.1,ylabel={绝对误差},restrict x to domain=0.1:1]
\addplot[violet,very thick] table[x=x,y=ee] {assets/tikz/F003_samples.dat};
\addlegendentry{$|e^x-R_E(x)|$}
\addplot[gray!80!black,thick,densely dotted] table[x=x,y=e4] {assets/tikz/F003_samples.dat};
\addlegendentry{$|e^x-T_4(x)|$}

\end{axis}
\end{tikzpicture}
\end{center}

上图将差值单独画出；$T_4=1+x+x^2/2+x^3/6+x^4/24$。
两种方法都匹配到四次项，但在所画区间内，分式的误差更小。
这只是当前函数与区间的比较，不能推广成“帕德永远优于泰勒”。

\newpage
\paragraph{二、对数：负数一侧更靠近定义域边界}
\begin{center}
\begin{tikzpicture}[font=\small]
\draw[black!30,rounded corners=2pt] (0,0) rectangle (14.4,8.2);
\node[anchor=west,font=\small\bfseries] at (0.5,7.72) {对数：曲线贴合，不代表误差为零};
\begin{axis}[width=11.2cm,height=6.1cm,at={(2.05cm,1.25cm)},anchor=south west,
axis lines=left,xlabel={$x$},tick label style={font=\small},label style={font=\small},
legend style={font=\small,draw=none,fill=white,at={(0.03,0.97)},anchor=north west},
legend cell align=left,unbounded coords=jump,xmin=-0.8,xmax=1,ymin=-1.9,ymax=1,ylabel={$y$},legend style={at={(0.97,0.05)},anchor=south east}]
\addplot[blue,very thick] table[x=x,y=log] {assets/tikz/F003_samples.dat};
\addlegendentry{$\ln(1+x)$}
\addplot[orange!85!black,thick,dashed] table[x=x,y=rl] {assets/tikz/F003_samples.dat};
\addlegendentry{$R_L(x)$}

\end{axis}
\end{tikzpicture}
\end{center}

零点两侧曲线都贴合，但向 $x=-1$ 靠近时要更谨慎；不能只看 $|x|$ 就保证精度。
\begin{center}
\begin{tikzpicture}[font=\small]
\draw[black!30,rounded corners=2pt] (0,0) rectangle (14.4,8.2);
\node[anchor=west,font=\small\bfseries] at (0.5,7.72) {对数：误差在零点两侧改变符号};
\begin{axis}[width=11.2cm,height=6.1cm,at={(2.05cm,1.25cm)},anchor=south west,
axis lines=left,xlabel={$x$},tick label style={font=\small},label style={font=\small},
legend style={font=\small,draw=none,fill=white,at={(0.03,0.97)},anchor=north west},
legend cell align=left,unbounded coords=jump,xmin=-0.5,xmax=0.5,ymin=-10,ymax=2,ylabel={$10^4[\ln(1+x)-R_L(x)]$},restrict x to domain=-0.5:0.5]
\addplot[violet,very thick] table[x=x,y expr={10000*\thisrow{el}}] {assets/tikz/F003_samples.dat};
\addplot[black!50,dashed,domain=-0.5:0.5]{0};
\end{axis}
\end{tikzpicture}
\end{center}

误差为负表示分式高于原函数，误差为正表示分式低于原函数。
图中纵坐标已乘 $10^4$；实际误差须再除以 $10^4$。

\newpage
\paragraph{三、三角函数：先保留奇偶结构，再讨论精度}
\begin{center}
\begin{tikzpicture}[font=\small]
\draw[black!30,rounded corners=2pt] (0,0) rectangle (14.4,8.2);
\node[anchor=west,font=\small\bfseries] at (0.5,7.72) {三角函数：同一尺度下看奇偶结构与局部贴合};
\node[anchor=west] at (0.5,6.95) {正弦：奇函数};
\begin{axis}[width=6.2cm,height=5.9cm,at={(0.5cm,0.65cm)},anchor=south west,axis lines=middle,xmin=-1.2,xmax=1.2,ymin=-1.3,ymax=1.3,xlabel={$x$},ylabel={$y$},xtick={-1,1},ytick={-1,1},legend style={font=\small,draw=none,at={(0.97,0.03)},anchor=south east}]
\addplot[blue,very thick] table[x=x,y=sin] {assets/tikz/F003_samples.dat};
\addlegendentry{原函数}
\addplot[orange!85!black,thick,dashed] table[x=x,y=rs] {assets/tikz/F003_samples.dat};
\addlegendentry{分式}
\end{axis}
\node[anchor=west] at (7.5,6.95) {余弦：偶函数};
\begin{axis}[width=6.2cm,height=5.9cm,at={(7.5cm,0.65cm)},anchor=south west,axis lines=middle,xmin=-1.2,xmax=1.2,ymin=-1.3,ymax=1.3,xlabel={$x$},ylabel={$y$},xtick={-1,1},ytick={-1,1},legend style={font=\small,draw=none,at={(0.97,0.03)},anchor=south east}]
\addplot[blue,very thick] table[x=x,y=cos] {assets/tikz/F003_samples.dat};
\addlegendentry{原函数}
\addplot[orange!85!black,thick,dashed] table[x=x,y=rc] {assets/tikz/F003_samples.dat};
\addlegendentry{分式}
\end{axis}

\end{tikzpicture}
\end{center}

$R_S$ 为奇函数、$R_C$ 为偶函数，分别保留正弦和余弦的对称性。
两幅图使用相同的坐标范围和尺度；奇偶性只是必要的结构检查，不能单独证明高阶精度。
\begin{center}
\begin{tikzpicture}[font=\small]
\draw[black!30,rounded corners=2pt] (0,0) rectangle (14.4,8.2);
\node[anchor=west,font=\small\bfseries] at (0.5,7.72) {正切：局部贴合须远离极点};
\begin{axis}[width=11.2cm,height=6.1cm,at={(2.05cm,1.25cm)},anchor=south west,
axis lines=left,xlabel={$x$},tick label style={font=\small},label style={font=\small},
legend style={font=\small,draw=none,fill=white,at={(0.03,0.97)},anchor=north west},
legend cell align=left,unbounded coords=jump,xmin=-1,xmax=1,ymin=-1.9,ymax=1.9,ylabel={$y$},restrict x to domain=-1:1,legend style={at={(0.03,0.97)},anchor=north west}]
\addplot[blue,very thick] table[x=x,y=tan] {assets/tikz/F003_samples.dat};
\addlegendentry{$\tan x$}
\addplot[orange!85!black,thick,dashed] table[x=x,y=rt] {assets/tikz/F003_samples.dat};
\addlegendentry{$R_T(x)$}

\end{axis}
\end{tikzpicture}
\end{center}

$R_T$ 也为奇函数。即使在这幅局部图中，接近左右边缘时仍可看到正切的差距变大。
\end{explanationbox}
